\chapter{Sorting and Randomized Algorithms}

\section{The Selection Problem}

\begin{definition}
\emph{Selection}: given an array $A$ of $n$ numbers and $0 \le k < n$, find the
item that would be at index $k$ if $A$ were sorted. Special cases: $k = 0$
(minimum), $k = n-1$ (maximum), $k = \floor{n/2}$ (median).
\end{definition}

Easy solutions: sort ($\Theta(n \log n)$), or use a heap ($\bigO(n + k\log n)$).
We want $\Theta(n)$.

\subsection{Partition}

Both QuickSelect and QuickSort rely on two subroutines:
\begin{itemize}
  \item \code{choose-pivot(A)}: return an index $p$ (simplest: the last index).
  \item \code{partition(A, p)}: rearrange $A$ so that all items $\le A[p]$ come
        first, then the pivot at index $i$, then all items $> A[p]$; return $i$.
\end{itemize}

An efficient in-place partition (Hoare-style, as in the course):
\begin{lstlisting}
partition(A, p):
  swap A[n-1] and A[p]          // pivot now at the end
  v <- A[n-1]
  i <- -1; j <- n-1
  loop
    do i <- i+1 while i < n and A[i] < v
    do j <- j-1 while j > 0 and A[j] > v
    if i >= j then break
    swap A[i] and A[j]
  swap A[n-1] and A[i]          // put pivot in place
  return i
\end{lstlisting}
Invariant: $A[0..i-1] \le v$ and $A[j+1..n-2] \ge v$. Runs in $\Theta(n)$ time
with $\bigO(1)$ extra space.

\subsection{QuickSelect}

\begin{lstlisting}
quick-select(A, k):
  p <- choose-pivot(A)
  i <- partition(A, p)
  if i = k then return A[i]
  else if i > k then return quick-select(A[0..i-1], k)
  else return quick-select(A[i+1..n-1], k - i - 1)
\end{lstlisting}

\paragraph{Worst case.} If the pivot is always the largest item and $k = 0$, we
recurse on size $n-1$ every time: $T(n) = T(n-1) + cn$, so $\Theta(n^2)$.

\paragraph{Best case.} The first pivot is the answer: $\Theta(n)$.

\paragraph{Average case.} Assume all $n!$ input permutations are equally likely
(equivalently, the pivot's final position $i$ is uniform on $\{0,\ldots,n-1\}$
and the subarrays are again random). Then
\[
T^{\text{avg}}(n) \le cn + \frac1n \sum_{i=0}^{n-1} \max\{T(i), T(n-i-1)\}.
\]

\begin{theorem}
$T^{\text{avg}}(n) \in \Theta(n)$ for QuickSelect.
\end{theorem}
\begin{proof}[Proof sketch]
Call the pivot \emph{good} if $n/4 \le i < 3n/4$; this happens with probability
$1/2$. After a good pivot the recursion has size at most $3n/4$; after a bad one,
at most $n$. Then
$T(n) \le cn + \tfrac12 T(3n/4) + \tfrac12 T(n)$, so $T(n) \le 2cn + T(3n/4)$,
which unrolls to $T(n) \le 2cn \sum_{j\ge0} (3/4)^j = 8cn$.
\end{proof}

\section{Randomized Algorithms}

The average-case analysis assumes the input is random, which we cannot control
(an adversary could feed sorted data). Instead we make the \emph{algorithm} random.

\begin{definition}
A \emph{randomized algorithm} uses random bits. For a fixed instance $I$ and
random outcome $R$, the running time is $T(I, R)$. The \emph{expected running
time} is $T^{\text{exp}}(I) = \mathbb E_R[T(I,R)] = \sum_R T(I,R)\Pr[R]$, and the
\emph{worst-case expected time} is $T^{\text{exp}}(n) = \max_{|I|=n} T^{\text{exp}}(I)$.
\end{definition}

Two approaches:
\begin{enumerate}
  \item Shuffle the input randomly first (in $\Theta(n)$ by Fisher--Yates), then run
        the deterministic algorithm.
  \item Choose the pivot uniformly at random: \code{p <- random(n)}.
\end{enumerate}
With a random pivot, \emph{every} input has the same expected behaviour as the
average-case analysis, so randomized QuickSelect has expected time $\Theta(n)$
on every instance. The worst case is still $\Theta(n^2)$ but occurs with tiny
probability and not for any particular input.

\begin{lstlisting}
shuffle(A):
  for i <- n-1 down to 1 do
    swap A[i] and A[random(i+1)]    // random(m) uniform in {0..m-1}
\end{lstlisting}

\begin{aside}
There is also a deterministic $\Theta(n)$ worst-case selection algorithm
(median-of-medians, choosing a pivot from medians of groups of 5). It is mostly
of theoretical interest because of its large constant; CS 341 covers it.
\end{aside}

\section{QuickSort}

\begin{lstlisting}
quick-sort(A):
  if n <= 1 then return
  p <- choose-pivot(A)
  i <- partition(A, p)
  quick-sort(A[0..i-1])
  quick-sort(A[i+1..n-1])
\end{lstlisting}

\begin{itemize}
  \item Worst case: $T(n) = T(n-1) + \Theta(n) \in \Theta(n^2)$ (e.g.\ sorted
        input with last-element pivot).
  \item Best case: pivot always the median: $T(n) = 2T(n/2) + \Theta(n) \in
        \Theta(n \log n)$.
\end{itemize}

\subsection{Average-Case Analysis}

Let $T(n)$ be the average number of comparisons. Each pivot position $i$ is
equally likely, so
\[
T(n) = (n - 1) + \frac1n \sum_{i=0}^{n-1} \big(T(i) + T(n-1-i)\big)
     = (n-1) + \frac2n \sum_{i=0}^{n-1} T(i).
\]

\begin{theorem}
The average (and randomized expected) running time of QuickSort is
$\Theta(n \log n)$; precisely, the expected number of comparisons is
$2(n+1)H_n - 4n \approx 1.39\, n \log_2 n$.
\end{theorem}

\begin{proof}[Proof sketch 1: the recurrence]
Multiply by $n$: $nT(n) = n(n-1) + 2\sum_{i<n} T(i)$. Subtract the same identity
for $n-1$: $nT(n) - (n-1)T(n-1) = 2(n-1) + 2T(n-1)$, so
$nT(n) = (n+1)T(n-1) + 2(n-1)$. Divide by $n(n+1)$:
\[
\frac{T(n)}{n+1} = \frac{T(n-1)}{n} + \frac{2(n-1)}{n(n+1)}
\le \frac{T(n-1)}{n} + \frac{2}{n+1}.
\]
Telescoping gives $T(n)/(n+1) \le 2 H_{n+1}$, so $T(n) \in \bigO(n\log n)$.
\end{proof}

\begin{proof}[Proof sketch 2: indicator variables]
Let $z_1 < \cdots < z_n$ be the items in sorted order. Two items are compared at
most once, and only if one of them is the pivot. Items $z_i$ and $z_j$ ($i<j$)
are compared iff the first pivot chosen from $\{z_i, \ldots, z_j\}$ is $z_i$ or
$z_j$, which has probability $\frac{2}{j - i + 1}$. By linearity of expectation,
\[
\mathbb E[\#\text{comparisons}] = \sum_{i<j} \frac{2}{j-i+1}
\le \sum_{i=1}^{n} \sum_{d=2}^{n} \frac{2}{d} \le 2n H_n \in \bigO(n \log n).
\]
\end{proof}

\subsection{Practical Improvements}
\begin{itemize}
  \item Random pivot, or median-of-three, to avoid sorted-input worst case.
  \item Recurse on the \emph{smaller} side first and loop on the larger side, so
        the stack depth is $\bigO(\log n)$ worst case.
  \item Stop recursing at small sizes (e.g.\ $n \le 10$) and finish with one
        pass of insertion sort.
  \item Handle many duplicates with three-way partition ($<$, $=$, $>$).
\end{itemize}

\section{Lower Bound for Comparison Sorting}

\begin{definition}
A \emph{comparison-based} sorting algorithm accesses keys only by comparing pairs
(\code{A[i] < A[j]?}).
\end{definition}

\begin{theorem}
Every comparison-based sorting algorithm requires $\Omega(n \log n)$ comparisons
in the worst case (and also on average).
\end{theorem}
\begin{proof}
Model any such algorithm on inputs of size $n$ as a \emph{decision tree}: each
internal node is a comparison, each branch an outcome, each leaf an output
permutation. Each of the $n!$ input orderings of distinct keys must reach a
different leaf (otherwise two different inputs get the same rearrangement, and
one of them is sorted wrongly). A binary tree with at least $n!$ leaves has
height at least $\log_2(n!)$. Finally
\[
\log(n!) = \sum_{i=1}^n \log i \ge \sum_{i=n/2}^{n} \log\frac n2
\ge \frac n2 \log \frac n2 \in \Omega(n \log n).
\]
The worst-case number of comparisons is the height, so it is $\Omega(n\log n)$.
\end{proof}

So MergeSort and HeapSort are asymptotically optimal among comparison sorts. To
beat $n \log n$ we must use more than comparisons, e.g.\ the digits of the keys.

\section{Non-Comparison Sorting}

Assume keys are integers in $\{0, \ldots, R-1\}$ (or strings of $m$ digits in
radix $R$).

\subsection{Counting Sort (Key-Indexed Sort)}

\begin{lstlisting}
counting-sort(A, R):           // keys in 0..R-1
  C[0..R-1] <- all 0
  for i <- 0 to n-1 do C[A[i]] <- C[A[i]] + 1
  I[0] <- 0                    // I[k] = first index for key k
  for k <- 1 to R-1 do I[k] <- I[k-1] + C[k-1]
  B <- copy of A
  for i <- 0 to n-1 do
    A[I[B[i]]] <- B[i]
    I[B[i]] <- I[B[i]] + 1
\end{lstlisting}
Time $\Theta(n + R)$, space $\Theta(n + R)$. It is \emph{stable}: equal keys
keep their relative order, because we scan $B$ left to right.

\subsection{LSD Radix Sort}

Keys have $m$ digits in base $R$. Sort by the least significant digit first,
using a stable sort (counting sort) at each step.
\begin{lstlisting}
LSD-radix-sort(A):
  for d <- m down to 1 do      // digit m is least significant
    counting-sort A by digit d
\end{lstlisting}
Correctness (by induction): after processing digits $d..m$, the array is sorted
by the last $m - d + 1$ digits; stability ensures ties on digit $d$ keep the
order given by the less significant digits. Time $\Theta(m(n + R))$, space
$\Theta(n + R)$.

\begin{example}
Sort $\{123, 230, 21, 320, 210, 232, 101\}$ ($R = 10$, $m = 3$, pad 21 as 021):
\begin{itemize}
  \item by last digit: 230, 320, 210, 021, 101, 232, 123;
  \item by middle digit: 101, 210, 320, 021, 123, 230, 232;
  \item by first digit: 021, 101, 123, 210, 230, 232, 320.
\end{itemize}
\end{example}

\subsection{MSD Radix Sort}

Sort by the most significant digit, then recursively sort each group by the
remaining digits.
\begin{lstlisting}
MSD-radix-sort(A, l, r, d):
  if l < r then
    counting-sort A[l..r] by digit d
    if d < m then
      for each digit value k with a nonempty group A[l_k..r_k] do
        MSD-radix-sort(A, l_k, r_k, d+1)
\end{lstlisting}
Time $\bigO(m n R)$ in the worst case due to many tiny subproblems that each pay
$\Theta(R)$; in practice one switches to another sort for small groups. MSD can
stop early on distinct prefixes, and handles variable-length strings naturally.

\section{Summary of Sorting Algorithms}

\begin{center}
\begin{tabular}{lllll}
\toprule
Algorithm & Worst & Average/Expected & Space & Stable?\\
\midrule
Selection sort & $\Theta(n^2)$ & $\Theta(n^2)$ & $\Theta(1)$ & no\\
Insertion sort & $\Theta(n^2)$ & $\Theta(n^2)$ (best $\Theta(n)$) & $\Theta(1)$ & yes\\
MergeSort & $\Theta(n\log n)$ & $\Theta(n\log n)$ & $\Theta(n)$ & yes\\
HeapSort & $\Theta(n\log n)$ & $\Theta(n\log n)$ & $\Theta(1)$ & no\\
QuickSort (random) & $\Theta(n^2)$ & $\Theta(n\log n)$ expected & $\Theta(\log n)$ & no\\
Counting sort & $\Theta(n+R)$ & $\Theta(n+R)$ & $\Theta(n+R)$ & yes\\
LSD radix sort & $\Theta(m(n+R))$ & $\Theta(m(n+R))$ & $\Theta(n+R)$ & yes\\
\bottomrule
\end{tabular}
\end{center}

\begin{itemize}
  \item QuickSelect solves selection in $\Theta(n)$ expected (randomized) time.
  \item Randomization moves the average case from ``random input'' to ``random
        choices'', giving guarantees for every input.
  \item Comparison sorts need $\Omega(n \log n)$ by the decision-tree argument;
        radix sorts sidestep this by looking at digits.
\end{itemize}
